% Vista científica exacta materializada desde una rama condicional.
% Fuente: source/demostraciones_primarias/14_05b_extension_superviviente_caracter.tex; rama: HMTIncludeCharacterBlock.
% No modifica el contenido matemático de la rama de origen.
\chapter{Caractère positif de la signature et composition de registres}
\section{Caractère formel universel}

Soit \(\mathcal P=\RR^5\), de coordonnées
\[
 X=(X_A,X_C,X_4,X_{120},X_{270}).
\]
Pour \(v=(n_A,s_C,k,\nu_{120},\nu_{270})\in\ZZ^5\), nous définissons
\[
 \ell_{\rm form}(v;X)
 =n_AX_A+s_CX_C+kX_4
  +\nu_{120}X_{120}+\nu_{270}X_{270},
\]
\[
 \chi_{\rm form}(v;X)=\exp\!\bigl(\ell_{\rm form}(v;X)\bigr).
\]
Ainsi, \(\chi_{\rm form}\) prend ses valeurs dans le groupe multiplicatif
\(\operatorname{Fun}(\mathcal P,\RR_{>0})\), muni du produit ponctuel. Il n'a pas encore
sélectionné de coefficients globaux ni d'espèce.

\begin{theorem}[Caractère universel de la signature]
\label{thm:caracter-formal-ruta}
L'application
\[
 \chi_{\rm form}:(\ZZ^5,+)\longrightarrow
 \bigl(\operatorname{Fun}(\mathcal P,\RR_{>0}),\cdot\bigr)
\]
est un homomorphisme. La composition
\[
 \chi_{\rm route}^{\rm form}
 :=\chi_{\rm form}\circ L_{\rm tick}
\]
évalue formellement tout registre enrichi et n'introduit pas d'unité
physique.
\end{theorem}

\begin{proof}
Pour tout \(X\in\mathcal P\), l'exposant
\(\ell_{\rm form}(\,\cdot\,;X)\) est additif sur \(\ZZ^5\). Par conséquent,
\[
\chi_{\rm form}(v+w;X)
=\chi_{\rm form}(v;X)\chi_{\rm form}(w;X),
\]
et l'identité vaut comme égalité de fonctions positives.
\end{proof}

Le caractère formel précède la construction des coefficients globaux.
Cette section ne définit pas encore de caractère numérique de masse ni de
carte d'unités.

\section{Section interne minimale sur les 81 cellules}

L'atlas fini contient 81 cellules orientées. Leur projection sur la carte
de famille de chemin donne 56 valeurs. Notons cette projection
\[
 \mathcal R:\mathcal C_9\longrightarrow\mathcal R_{56}.
\]
L'histogramme exact de ses fibres est
\[
 41\times1,
 \qquad10\times2,
 \qquad2\times3,
 \qquad1\times4,
 \qquad2\times5,
\]
et leur somme est 81.

L'origine et l'ordre lexicographique de la carte étant fixés, soit
\(\kappa(x)\in\{0,1,2,3,4\}\) le rang de \(x\) au sein de sa fibre.

\begin{theorem}[Sélecteur interne minimal]
\label{thm:selector-interno-minimo}
L'application
\[
 x\longmapsto\bigl(\mathcal R(x),\kappa(x)\bigr)
\]
est injective sur les 81 cellules. Aucun alphabet de mémoire comportant moins de
cinq symboles ne peut relever \(\mathcal R\) injectivement sur l'ensemble de
l'atlas.
\end{theorem}

\begin{proof}
Au sein de chaque fibre, \(\kappa\) attribue des rangs distincts ; entre les fibres, la
première coordonnée est déjà distincte. Il existe une fibre de taille cinq, de sorte que
le principe des tiroirs impose à tout relèvement injectif de disposer d'au moins
cinq valeurs de mémoire. Le rang \(0,\ldots,4\) atteint cette borne.
\end{proof}

La canonicité est relative à l'origine et à l'ordre déclarés. \(\kappa\) n'est
ni une charge, ni une saveur, ni une masse. Les treize classes grossières de l'atlas produisent
des multisections finies ; seule la classe centrale possède le point fixe \((5,5)\)
sous l'inversion centrale. Les douze autres nécessitent une information orientée
pour sélectionner un point et n'admettent pas de sélecteur ponctuel équivariant à partir
de la seule classe grossière.

\subsection{Inversion de Möbius et caractères d'aller et de retour}

Sur le mot de signes dodécaphasique, nous définissons
\[
 C_M(\tau)_m=-\tau_{11-m}.
\]
Cette opération combine l'inversion de l'ordre de lecture avec le changement de
feuillet. Un caractère linéaire
\(\ell_a(\tau)=\sum_ma_m\tau_m\), avec
\(a_m=a_{11-m}\), satisfait
\[
 \ell_a(C_M\tau)=-\ell_a(\tau),
\]
tandis que toute forme quadratique symétrique compatible avec la réversion
satisfait
\[
 Q(C_M\tau)=Q(\tau).
\]
Ainsi, les coordonnées linéaires distinguent l'orientation d'aller et de retour et les
quadratiques conservent le contenu de corrélation. Cette parité explique
pourquoi \(s_C\) change de signe tandis que \(\nu_{270}\) peut rester
invariant. Dans l'atlas, les sept formes de mot ont pour multiplicités
\[
 54,2,2,2,7,7,7;
\]
c'est pourquoi l'antifeuillet se réalise comme correspondance entre fibres et non comme
une involution bijective obtenue en ignorant leurs multiplicités.

\section{Composition orientée de registres avant la projection}

Soient \(\Gamma\) et \(\Lambda\) deux histoires enrichies, \(B\) une
sous-histoire commune et \(g\) le transport qui identifie phase, feuillet,
orientation et position à l'interface. Au niveau du registre, on définit
\[
 \boxed{
 \Gamma\star_B\Lambda
 =\bigl(n_\Gamma+n_\Lambda-n_B,
        e_\Gamma+g e_\Lambda-e_B\bigr).}
\]
La définition présuppose aussi l'identification du bord, des
prédécesseurs torsionnels, de la couronne et de la mémoire antifeuillet.

\begin{proposition}[Obstruction à la descente pentadimensionnelle de la composition de registres]
\label{prop:no-descenso-empalme-cinco}
En général, il n'existe pas d'opération binaire sur \(\ZZ^5\) qui récupère
\(L_{\rm tick}(\Gamma\star_B\Lambda)\) à partir de
\(L_{\rm tick}(\Gamma)\) et \(L_{\rm tick}(\Lambda)\) seulement.
\end{proposition}

\begin{proof}
L'opération compose les douze événements avant d'appliquer
\(\operatorname{sgn}\) et les autocorrélations. Deux histoires peuvent avoir
le même quintuplet et des valeurs distinctes de \(M_5\) ou \(A_6\). En les composant
avec une troisième histoire, une somme
\(e_a+e_{a+4}+e_{a+8}\) peut franchir zéro dans l'une et non dans l'autre, modifiant
\(\nu_{120}^{\rm ev}\). Les losanges finis fournissent des témoins explicites.
Par conséquent, la projection sur cinq entiers n'est pas
une congruence pour \(\star_B\).
\end{proof}

La proposition formalise une règle opératoire essentielle : les extensions se
composent ; les signatures se lisent ensuite. La particule composée ne
s'obtient pas non plus en additionnant aveuglément deux lignes d'un tableau, mais en composant leurs mémoires
compatibles et en évaluant le résultat.

\section{Dépendance dodécaphasique du caractère de persistance}

Les moteurs réduits de période 54 imposent
\[
 e_{m+6}=e_m.
\]
Il en découle les congruences
\[
 n_A\equiv0\pmod2,
 \qquad
 s_C\equiv0\pmod2,
 \qquad
 \nu_{270}\equiv0\pmod4.
\]
Les signatures de la reconstruction scalaire abrégée violent au moins une de
ces congruences. Leur généalogie exige donc le registre enrichi et
ne procède pas de ce moteur réduit avec les règles déclarées.

Ce résultat n'est pas une obstruction à la dynamique TPK enrichie. Il démontre
quelle mémoire elle nécessite : une dodécaphase véritable, avec éventuellement
\(A_6\ne0\), une composition orientée et des prédécesseurs torsionnels conservés. L'ablation
évite d'attribuer à TPK une limitation qui appartient à une réduction
périodique concrète.

\section{Factorisation par le quotient de registre et réalisation physique}

Dans le cadre de cette section, on déclare la relation
\[
\gamma\sim_{\mathcal L}\gamma'
\quad\Longleftrightarrow\quad
\text{\(\gamma,\gamma'\) présentent la même extension et }
\mathcal L(\gamma)=\mathcal L(\gamma').
\]
Nous notons
\(\mathcal E_{108}^{\rm surv}:=
\Gamma_{108}^{\rm enr}/\!\sim_{\mathcal L}\)
le quotient relatif à cette carte. Soit
\[
 \pi_{\rm ext}:\Gamma_{108}^{\rm enr}\longrightarrow
 \mathcal E_{108}^{\rm surv}
\]
la projection canonique. Par construction, l'extracteur est constant sur ses
fibres. L'affirmation plus forte selon laquelle toute carte admissible d'une
même extension produit automatiquement le même registre exigerait de démontrer
la compatibilité entre cartes et n'est pas présupposée ici.

\begin{theorem}[Factorisation relative par le quotient de registre]
\label{thm:descenso-caracter-extension}
Il existe une unique application
\[
 \overline L:\mathcal E_{108}^{\rm surv}\longrightarrow\ZZ^5
\]
tel que
\[
 \boxed{L_{\rm tick}=\overline L\circ\pi_{\rm ext}.}
\]
Par conséquent, \(\chi_{\rm form}\circ\overline L\) est une fonction de la classe dans
le quotient relatif et non du représentant de chemin choisi pour la lire.
\end{theorem}

\begin{proof}
Si \(\pi_{\rm ext}(\gamma)=\pi_{\rm ext}(\gamma')\), les deux présentations
possèdent le même registre par la relation déclarée et, d'après les formules
précédentes, le même quintuplet. Nous définissons
\(\overline L([\gamma])=L_{\rm tick}(\gamma)\). La constance prouve que la
définition est indépendante du représentant ; la surjectivité de l'application
au quotient prouve l'unicité.
\end{proof}

La construction interne atteint ainsi
\[
 \Omega_{\HMT}^{\rm surv}
 \longrightarrow\mathcal E_{108}^{\rm surv}
 \xrightarrow{\ \overline L\ }\ZZ^5
 \xrightarrow{\ \chi_{\rm form}\ }
 \operatorname{Fun}(\mathcal P,\RR_{>0}).
\]
Le contenu du théorème est positif et complet dans ce domaine : une
extension survivante détermine sa classe de registre, la classe détermine
une unique signature et la signature détermine le caractère formel. La spécialisation
globale et la hiérarchie complète des tours sont construites dans les chapitres
suivants ; les
comparaisons avec les noms et les unités physiques sont présentées ensuite, sans
modifier cette chaîne.
